% ============================================================================
\bigskip
\appendix
\renewcommand{\thesection}{\Alph{section}}
\titleformat{\section}{\normalfont\Large\bfseries}{Phụ lục \thesection.}{0.6em}{}
% ============================================================================

\section{Hướng dẫn tái lập}
\label{app:reproduce}

Các lệnh dưới đây chạy từ thư mục gốc của kho mã \code{uit-vietnamese-mrc}. Mã nguồn nằm trong
\code{src/} nhưng không được cài như một gói, nên \textbf{mọi script cần biến môi trường
\code{PYTHONPATH=src}}; thiếu biến này, các script dừng với lỗi \code{ModuleNotFoundError}.

\begin{lstlisting}[language=bash]
# 0. Moi truong (Python >= 3.12; tren Apple Silicon PyTorch tu dung MPS)
uv venv --python 3.12 .venv && source .venv/bin/activate
uv pip install -r requirements.txt
export PYTHONPATH=src

# 1. Tai du lieu -> data/raw/viquad2_{train,validation,test}.json
python scripts/fetch_data.py

# 2. Kiem thu nhanh (khong tai mo hinh, ~2 giay) va day du
pytest -m "not slow"
pytest --cov=mrc

# 3. Fine-tune tren toan bo train (~35 phut/epoch voi mBERT tren M5 Pro)
python scripts/finetune.py --model bert-base-multilingual-cased --out models/mbert
python scripts/finetune.py --model uitnlp/visobert --out models/visobert \
       --epochs 3 --lr 5e-5 --max-answer-len 64

# 4. Danh gia tren cung mau 500 cau (seed 42); --full de chay ca 3.814 cau
python scripts/run_eval.py --models baseline xlmr mbert visobert --limit 500

# 5. Doi chieu gia thuyet va sinh hinh tu results/*.json
python scripts/check_hypotheses.py
python scripts/make_figures.py

# 6. Ung dung web
streamlit run app/streamlit_app.py
\end{lstlisting}

\paragraph{Tài liệu phân tích đi kèm.} Bản phân tích kết quả đầy đủ --- gồm bảng đối chiếu từng
giả thuyết đăng ký trước, chẩn đoán lần huấn luyện ViSoBERT, phân rã chênh lệch theo nhóm câu,
và toàn bộ bằng chứng kiểm định tính tái lập ở Phụ lục~\ref{app:repro-evidence} --- nằm ở
\code{08\_PHAN\_TICH/PHAN\_TICH\_KET\_QUA.md} trong thư mục nộp. Tài liệu đó dẫn nguồn tới từng
tệp trong \code{results/} cho mỗi con số, và là bản ghi chi tiết mà Chương~\ref{sec:experiments}
tóm lược.

Kết quả đánh giá được ghi vào \code{results/eval\_<mô hình>\_validation.json}; đường cong huấn luyện
vào \code{results/training\_curve\_<mô hình>.json}; hình vẽ vào \code{results/figures/}. Các thống kê
và hình bổ sung chỉ có trong báo cáo này (Bảng~\ref{tab:data-stats}, \ref{tab:wh},
\ref{tab:decomp}, \ref{tab:tokenization}, các khoảng tin cậy, Hình~\ref{fig:dataset-composition},
\ref{fig:length-dist}) được sinh bằng các script trong thư mục \code{evidence/} đi kèm báo cáo.

\section{Bằng chứng tái lập và toàn vẹn gói nộp}
\label{app:repro-evidence}

Phụ lục này ghi lại bằng chứng \emph{đo được} cho các khẳng định về chất lượng kỹ thuật ở
Mục~\ref{sec:testing}, thay vì khẳng định suông. Phụ lục~\ref{app:reproduce} nói \emph{cách} chạy
lại; phụ lục này nói kết quả khi đã chạy lại thật. Mọi phép đo dưới đây được thực hiện trên gói nộp
thật, ở commit \code{c9a6de9}.

\subsection{Bộ kiểm thử chạy trong container}

\begin{table}[htbp]
\centering
\caption{Kết quả kiểm thử và khởi động, đo trên gói nộp}
\label{tab:repro-tests}
\renewcommand{\arraystretch}{1.2}
\small
\begin{tabularx}{\textwidth}{@{}L L l@{}}
\toprule
\textbf{Phép kiểm} & \textbf{Lệnh} & \textbf{Kết quả} \\
\midrule
Dựng ảnh & \code{docker compose build app} & \yes\ sạch, ảnh 2,49 GB \\
Test nhanh & \code{docker compose run -{}-rm tests} & \yes\ \textbf{423 đạt}, 51 bỏ chọn, \textbf{0 bị bỏ qua} \\
Test đầy đủ & \code{docker compose run -{}-rm tests-full} & \yes\ \textbf{474 đạt}, 0 hỏng, 0 bị bỏ qua \\
Khởi động demo & \code{docker compose up -{}-wait app} & \yes\ healthy; cả 6 màn hình trả HTTP 200 \\
\bottomrule
\end{tabularx}
\end{table}

\noindent Chi tiết đáng chú ý là cột cuối: \textbf{không có test nào bị bỏ qua}. Hai test cấm gõ
tay số đo (Mục~\ref{sec:testing}) đọc thư mục \code{report/} và hai tệp \code{slides/index.html},
\code{slides/build\_deck.js} từ
\emph{bên trong} container; nếu \code{.dockerignore} loại nhầm hai thư mục đó thì hai test này sẽ
\emph{lặng lẽ bỏ qua} thay vì báo đỏ. Con số 0 bị bỏ qua vì vậy là bằng chứng rằng chúng đã thật
sự chạy --- một test bị bỏ qua âm thầm không bảo vệ được gì.

\subsection{Mô phỏng máy người chấm}

Phép kiểm nghiêm nhất là xoá hẳn ảnh Docker khỏi máy rồi chạy gói đúng như người chấm sẽ chạy:

\begin{lstlisting}[language=bash]
docker image rm vietnamese-mrc:latest    # -> Deleted
./run.sh test                            # -> Loaded image: vietnamese-mrc:latest
                                         # -> 423 passed, 51 deselected
\end{lstlisting}

\noindent Gói \textbf{tự đứng được}: không cần mạng, không cần dựng lại ảnh, không cần Python
trên máy chủ.

\subsection{Demo và báo cáo có khớp nhau không}

Đây là bất biến số một của dự án (Mục~\ref{sec:testing}): con số trên màn hình và con số trong báo
cáo phải đến từ cùng một nguồn. Phép kiểm được thực hiện bằng suy luận thật trong container, chứ
không bằng cách đọc lại tệp kết quả.

\begin{table}[htbp]
\centering
\caption{Đối chiếu demo với báo cáo, và tính đơn điệu của ngưỡng từ chối $\tau$}
\label{tab:repro-demo}
\renewcommand{\arraystretch}{1.2}
\small
\begin{tabular}{@{}l r@{}}
\toprule
\textbf{Phép đo} & \textbf{Kết quả} \\
\midrule
Tỉ lệ từ chối đo trực tiếp trong container (mBERT, $\tau = 0{,}0$) & 40\% (2/5 câu) \\
\code{impossible\_EM} công bố ở Bảng~\ref{tab:ans-imp} & 41,01 \\
\bottomrule
\end{tabular}
\end{table}

\noindent Hai con số khớp nhau trong phạm vi cỡ mẫu của phép đo tại chỗ (5 câu): demo và báo cáo
đang đọc cùng một thực tại. Quét tiếp ngưỡng $\tau$ trên chính 5 câu đó --- đúng thao tác mà tài
liệu hướng dẫn đề nghị người chấm tự làm --- cho số câu bị từ chối là 2, 2, 3, 3, 5 tại
$\tau = 0{,}0$; $0{,}5$; $1{,}0$; $1{,}5$; $3{,}0$: \textbf{đơn điệu tăng}, đúng như thiết kế ở
Mục~\ref{sec:inference}. Biên độ \code{null\_delta} giữ nguyên giá trị $-2{,}371$ ở cả
$\tau = 0{,}0$ lẫn $\tau = 1{,}5$, xác nhận rằng ngưỡng chỉ dịch ranh giới quyết định chứ không
bóp méo điểm số của mô hình.

\subsection{Tính toàn vẹn của gói nộp và của dữ liệu}

\begin{table}[htbp]
\centering
\caption{Kiểm tra tính toàn vẹn của gói nộp và của các split dữ liệu}
\label{tab:repro-integrity}
\renewcommand{\arraystretch}{1.2}
\small
\begin{tabularx}{\textwidth}{@{}L L@{}}
\toprule
\textbf{Hạng mục} & \textbf{Kết quả} \\
\midrule
Kích thước gói demo & 1,5 GB (ảnh nén 502 MB) \\
Commit ghim trong \code{MANIFEST.txt} & \code{c9a6de9} \\
Đối chiếu manifest & \yes\ 137 tệp liệt kê $=$ 137 tệp thực tế \\
Kiểm tra nén & \yes\ \code{gzip -t} \\
Rác lọt vào gói & \yes\ không (\code{.DS\_Store}, \code{\_\_pycache\_\_}, \code{.venv}, \code{.omc} đều sạch)$^\ast$ \\
\midrule
Đoạn văn dùng chung train $\cap$ validation & \yes\ \textbf{0} --- không có rò rỉ \\
Đoạn văn dùng chung train $\cap$ test & 449 (không ảnh hưởng --- test là tập ẩn nhãn) \\
Câu chấm được trong test split & 0 --- toàn bộ đáp án vàng rỗng (Mục~\ref{sec:blind}) \\
\bottomrule
\end{tabularx}
\par\vspace{2pt}
{\footnotesize $^\ast$ Phép đối chiếu manifest đáng chạy lại \emph{ngay trước khi nộp}: chỉ cần mở một
terminal có công cụ ghi trạng thái vào thư mục gói là số tệp thực tế lệch khỏi manifest, dù nội
dung nộp không đổi. Lệnh kiểm: so \code{find . -type f} với danh sách trong \code{MANIFEST.txt}.}
\end{table}

\noindent Dòng quan trọng nhất là \textbf{0 đoạn văn dùng chung giữa train và validation}. Đây là
bằng chứng trực tiếp rằng mọi điểm số trong chương này không bị thổi phồng do rò rỉ theo nhóm
(Mục~\ref{sec:datalayer}) --- và cũng là lý do mọi kết quả được báo cáo trên validation chứ không
phải trên test.

\section{Cấu trúc mã nguồn}
\label{app:structure}

\begin{lstlisting}
uit-vietnamese-mrc/
|-- src/
|   |-- mrc/                  # thu vien loi, co kiem thu
|   |   |-- normalize.py      #   chuan hoa chuoi (quyet dinh A/B/C)
|   |   |-- metrics.py        #   EM, token-F1, tong hop SQuAD-2.0
|   |   |-- data.py           #   Example, chong ro ri, assert_gradeable, lay mau con
|   |   |-- tagging.py        #   nhom do dai, loai cau hoi (heuristic)
|   |   |-- predictor.py      #   giao dien Predictor + do do tre
|   |   |-- baseline_tfidf.py #   baseline TF-IDF + cosine
|   |   |-- windowing.py      #   cua so truot tu cai, chon span, offset
|   |   |-- transformer_qa.py #   suy luan voi lop QA
|   |   |-- features.py       #   gan nhan token cho huan luyen
|   |   |-- training.py       #   lich hoc, chon epoch, chan doan duong cong
|   |   |-- evaluate.py       #   harness danh gia + provenance
|   |   `-- device.py         #   chon MPS / CUDA / CPU
|   |-- evaluation/           # logic dong lenh cua run_eval
|   |-- reporting/            # sinh hinh ve tu results/*.json
|   `-- demo/logic.py         # logic ung dung web, tach khoi Streamlit
|-- scripts/                  # vo mong: fetch_data, finetune, run_eval,
|                             #          make_figures, check_hypotheses
|-- app/streamlit_app.py      # giao dien web
|-- tests/                    # 474 test (51 danh dau slow)
|-- results/                  # ket qua JSON, nhat ky, gia thuyet, hinh ve
|-- docs/                     # SOLUTION.md, PLAN_TDD.md
|-- requirements.txt          # phien ban thu vien duoc ghim
`-- pyproject.toml
\end{lstlisting}

\section{Bảng truy vết số liệu}
\label{app:trace}

Bảng~\ref{tab:trace} ánh xạ các nhóm số liệu trong báo cáo về nguồn gốc của chúng, để người đọc
có thể kiểm tra lại bất kỳ con số nào.

\begin{table}[H]
\centering
\caption{Nguồn gốc các nhóm số liệu trong báo cáo}
\label{tab:trace}
\renewcommand{\arraystretch}{1.2}
\small
\begin{tabularx}{\textwidth}{@{}L L@{}}
\toprule
\textbf{Số liệu} & \textbf{Nguồn} \\
\midrule
EM, F1, tách nhóm, độ trễ (Bảng~\ref{tab:main}, \ref{tab:ans-imp}, \ref{tab:qtype}, \ref{tab:length}, \ref{tab:cost}) & \code{results/eval\_*\_validation.json} \\
Dự đoán mẫu (Bảng~\ref{tab:qualitative}) & Trường \code{sample\_predictions} của các tệp trên \\
Đường cong huấn luyện, siêu tham số (Bảng~\ref{tab:curves}, \ref{tab:hparams}) & \code{results/training\_curve\_*.json}, \code{results/finetune\_*.log} \\
Khoảng tin cậy, phân rã chênh lệch (Bảng~\ref{tab:decomp}) & \code{evidence/ci.py} $\to$ \code{evidence/confidence\_intervals.txt} \\
Thống kê dữ liệu, từ để hỏi (Bảng~\ref{tab:data-stats}, \ref{tab:wh}) & \code{evidence/stats.py} $\to$ \code{evidence/data\_stats.json} \\
Hình dữ liệu (Hình~\ref{fig:dataset-composition}, \ref{fig:length-dist}) & \code{evidence/figs.py} \\
Giả thuyết (Bảng~\ref{tab:hyp-quant}) & \code{results/hypotheses.md}, \code{results/hypotheses.json}, \code{scripts/check\_hypotheses.py} \\
Chẩn đoán ViSoBERT (Bảng~\ref{tab:confounds}, \ref{tab:threshold}) & Nhật ký commit \code{abb6d8f}, \code{c63b28c}, \code{388b0ba} \\
Phân mảnh token, số tham số (Bảng~\ref{tab:tokenization}) & \code{evidence/tokenizer\_stats.txt}; tokenizer, \code{config.json} và tệp trọng số trong \code{models/} \\
Số test, độ phủ (Bảng~\ref{tab:modules}) & \code{pytest} tại commit \code{c5b06ba} \\
Hình kết quả (Hình~\ref{fig:model-comparison}, \ref{fig:ans-imp}, \ref{fig:qtype}, \ref{fig:length}, \ref{fig:curves}) & \code{scripts/make\_figures.py} $\to$ \code{results/figures/} \\
Ảnh giao diện (Hình~\ref{fig:demo}) & Chụp từ ứng dụng chạy cục bộ với \code{models/mbert} \\
\bottomrule
\end{tabularx}
\end{table}
